%=====================================================================
%  Reference list, formatted to match the published article (APA-style,
%  the house style of Transportation Research Part F).
%
%  natbib label syntax:  \bibitem[Short(Year)Long]{key}
%    \citep{key}  -> "(Short, Year)"      e.g. "(Buehler & Pucher, 2021)"
%    \citet{key}  -> "Short (Year)"
%    \citet*{key} -> "Long (Year)"        e.g. "Buehler and Pucher (2021)"
%  The short form carries "&" (used inside parentheses) and the long form
%  carries "and" (used in running text), exactly as in the article. Use
%  \citet* for narrative citations of two-author works.
%
%  refs.bib carries the same entries in BibTeX form.
%=====================================================================

\begin{thebibliography}{99}
\setlength{\itemsep}{1pt plus 0.2pt}

\bibitem[Andersen et al.(2012)]{andersen2012collection}
Andersen, T., Breda, F., Weinreich, M., Jensen, N., Riisgaard-Dam, M., \&
Nielsen, M. K. (2012). \emph{Collection of cycle concepts 2012}. Technical
Report. Cycling Embassy of Denmark.
\url{https://bicycleinfrastructuremanuals.com/manuals1/Collection-of-Cycle-Concepts-2012_Denmark.pdf}.

\bibitem[Ballo et al.(2023)]{ballo2023ebike}
Ballo, L., de Freitas, L. M., Meister, A., \& Axhausen, K. W. (2023). The E-Bike
city as a radical shift toward zero-emission transport: Sustainable? Equitable?
Desirable? \emph{Journal of Transport Geography, 111}, Article 103663.
\url{https://doi.org/10.1016/j.jtrangeo.2023.103663}.
\url{https://www.sciencedirect.com/science/article/pii/S0966692323001357}.

\bibitem[Batista et al.(2023)]{batista2023exploring}
Batista, M., Berghoefer, F. L., \& Friedrich, B. (2023). Exploring pedestrian
and cyclist preferences for shared space design: A video-based online survey.
\emph{Transportation Research Interdisciplinary Perspectives, 22}, Article
100976. \url{https://doi.org/10.1016/j.trip.2023.100976}.

\bibitem[Beitel et al.(2018)]{beitel2018assessing}
Beitel, D., Stipancic, J., Manaugh, K., \& Miranda-Moreno, L. (2018). Assessing
safety of shared space using cyclist-pedestrian interactions and automated video
conflict analysis. \emph{Transportation Research. Part D, Transport and
Environment, 65}, 710--724. \url{https://doi.org/10.1016/j.trd.2018.10.001}.
\url{https://linkinghub.elsevier.com/retrieve/pii/S1361920917301694}.

\bibitem[Berghoefer \& Vollrath(2023)Berghoefer and Vollrath]{berghoefer2023prefer}
Berghoefer, F. L., \& Vollrath, M. (2023). Prefer what you like? Evaluation and
preference of cycling infrastructures in a bicycle simulator. \emph{Journal of
Safety Research}. \url{https://doi.org/10.1016/j.jsr.2023.09.013}.
\url{https://www.sciencedirect.com/science/article/pii/S0022437523001421}.

\bibitem[Beura et al.(2021)]{beura2021unsignalized}
Beura, S. K., Srivastava, A., \& Bhuyan, P. K. (2021). Unsignalized intersection
level of service: A bicyclist's perspective. \emph{International Journal of
Intelligent Transport Systems Research, 19}, 405--416.
\url{https://doi.org/10.1007/s13177-020-00244-z}.
\url{https://link.springer.com/10.1007/s13177-020-00244-z}.

\bibitem[Bogacz et al.(2020)]{bogacz2020comparison}
Bogacz, M., Hess, S., Calastri, C., Choudhury, C. F., Erath, A., van Eggermond,
M. A. B., Mushtaq, F., Nazemi, M., \& Awais, M. (2020). Comparison of cycling
behavior between keyboard-controlled and instrumented bicycle experiments in
virtual reality. \emph{Transportation Research Record, 2674}, 244--257.
\url{https://doi.org/10.1177/0361198120921850}. Publisher: SAGE Publications
Inc.

\bibitem[Boufous et al.(2018)]{boufous2018impact}
Boufous, S., Hatfield, J., \& Grzebieta, R. (2018). The impact of environmental
factors on cycling speed on shared paths. \emph{Accident Analysis and
Prevention, 110}, 171--176. \url{https://doi.org/10.1016/j.aap.2017.09.017}.
\url{https://www.sciencedirect.com/science/article/pii/S000145751730338X}.

\bibitem[Buehler \& Pucher(2008)Buehler and Pucher]{buehler2008making}
Buehler, R., \& Pucher, J. (2008). Making cycling irresistible: Lessons from the
Netherlands, Denmark and Germany. \emph{Transport Reviews, 28}, 495--528.
\url{https://doi.org/10.1080/01441640701806612}.

\bibitem[Buehler \& Pucher(2021)Buehler and Pucher]{buehler2021covid}
Buehler, R., \& Pucher, J. (2021). COVID-19 impacts on cycling, 2019--2020.
\emph{Transport Reviews, 41}, 393--400.
\url{https://doi.org/10.1080/01441647.2021.1914900}.

\bibitem[Che et al.(2021)]{che2021interaction}
Che, M., Wong, Y. D., Lum, K. M., \& Wang, X. (2021). Interaction behaviour of
active mobility users in shared space. \emph{Transportation Research. Part A,
Policy and Practice, 153}, 52--65.
\url{https://doi.org/10.1016/j.tra.2021.08.017}.
\url{https://linkinghub.elsevier.com/retrieve/pii/S0965856421002238}. Publisher:
Elsevier BV.

\bibitem[Chou et al.(2024)]{chou2024comparative}
Chou, K. Y., Paulsen, M., Jensen, A. F., Rasmussen, T. K., \& Nielsen, O. A.
(2024). Comparative modeling of risk factors for near-crashes from crowdsourced
bicycle airbag helmet data and crashes from conventional police data.
\emph{Journal of Safety Research, 91}, 465--480.
\url{https://doi.org/10.1016/j.jsr.2024.10.003}.
\url{https://linkinghub.elsevier.com/retrieve/pii/S0022437524001439}.

\bibitem[Cobb et al.(2021)]{cobb2021bicyclists}
Cobb, D. P., Jashami, H., \& Hurwitz, D. S. (2021). Bicyclists' behavioral and
physiological responses to varying roadway conditions and bicycle
infrastructure. \emph{Transportation Research. Part F, Traffic Psychology and
Behaviour, 80}, 172--188. \url{https://doi.org/10.1016/j.trf.2021.04.004}.
\url{https://www.sciencedirect.com/science/article/pii/S1369847821000851}.

\bibitem[Commission(2024)]{ec2024speed}
Commission, E. (2024). \emph{Current speed limit policies -- European
commission}.
\url{https://road-safety.transport.ec.europa.eu/eu-road-safety-policy/priorities/safe-road-use/safe-speed/archive/current-speed-limit-policies_en}.

\bibitem[Cyclenation(2014)]{cyclenation2014}
Cyclenation (2014). \emph{Making space for cycling -- a guide for new
developments and street renewals}. Technical Report.
\url{https://www.makingspaceforcycling.org/}.

\bibitem[Damien Masson(2024)]{masson2024latin}
Damien Masson (2024). \emph{Balanced Latin square online generator}.
\url{https://damienmasson.com/tools/latin_square/}.

\bibitem[Essa et al.(2018)]{essa2018road}
Essa, M., Hussein, M., \& Sayed, T. (2018). Road users' behavior and safety
analysis of pedestrian--bike shared space: Case study of Robson street in
Vancouver. \emph{Canadian Journal of Civil Engineering, 45}, 1053--1064.
\url{https://doi.org/10.1139/cjce-2017-0683}.
\url{http://www.nrcresearchpress.com/doi/10.1139/cjce-2017-0683}.

\bibitem[FGSV(2010)]{fgsv2010}
FGSV (2010). \emph{Empfehlungen für Radverkehrsanlagen}. Technical Report.
Forschungsgesellschaft für Straßen- und Verkehrswesen.
\url{https://www.fgsv-verlag.de/era}.

\bibitem[Friel et al.(2023)]{friel2023cyclists}
Friel, D., Wachholz, S., Werner, T., Zimmermann, L., Schwedes, O., \& Stark, R.
(2023). Cyclists' perceived safety on intersections and roundabouts -- a
qualitative bicycle simulator study. \emph{Journal of Safety Research, 87},
143--156. \url{https://doi.org/10.1016/j.jsr.2023.09.012}.
\url{https://www.sciencedirect.com/science/article/pii/S002243752300141X}.

\bibitem[Gkekas et al.(2020)]{gkekas2020perceived}
Gkekas, F., Bigazzi, A., \& Gill, G. (2020). Perceived safety and experienced
incidents between pedestrians and cyclists in a high-volume non-motorized shared
space. \emph{Transportation Research Interdisciplinary Perspectives, 4}, Article
100094. \url{https://doi.org/10.1016/j.trip.2020.100094}.
\url{https://www.sciencedirect.com/science/article/pii/S2590198220300051}.

\bibitem[Guo et al.(2023)]{guo2023psycho}
Guo, X., Tavakoli, A., Angulo, A., Robartes, E., Chen, T. D., \& Heydarian, A.
(2023). Psycho-physiological measures on a bicycle simulator in immersive
virtual environments: How protected/curbside bike lanes may improve perceived
safety. \emph{Transportation Research. Part F, Traffic Psychology and Behaviour,
92}, 317--336. \url{https://doi.org/10.1016/j.trf.2022.11.015}.
\url{https://www.sciencedirect.com/science/article/pii/S1369847822002832}.

\bibitem[Hamilton-Baillie(2008)]{hamiltonbaillie2008shared}
Hamilton-Baillie, B. (2008). Shared space: Reconciling people, places and
traffic. \emph{Built Environment, 34}, 161--181.
\url{https://doi.org/10.2148/benv.34.2.161}.

\bibitem[Huang et al.(2021)]{huang2021investigating}
Huang, C., Wang, L., Yu, H., Li, H., Zhang, J., Song, W., Lo, S., \& Rafaqat, W.
(2021). Investigating the influence of a cyclist on crowd behaviors on a shared
road. \emph{Journal of Statistical Mechanics: Theory and Experiment, 2021},
Article 083402. \url{https://doi.org/10.1088/1742-5468/ac0edd}.
\url{https://dx.doi.org/10.1088/1742-5468/ac0edd}. Publisher: IOP Publishing and
SISSA.

\bibitem[Hummer et al.(2005)]{hummer2005user}
Hummer, J. E., Rouphail, N., Hughes, R. G., Fain, S. J., Toole, J. L., Patten,
R. S., Schneider, R. J., Monahan, J. F., \& Do, A. (2005). User perceptions of
the quality of service on shared paths. \emph{Transportation Research Record,
1939}, 28--36. \url{https://doi.org/10.1177/0361198105193900104}. Publisher:
SAGE Publications Inc.

\bibitem[Jensen(2012)]{jensen2012pedestrian}
Jensen, S. U. (2012). \emph{Pedestrian and bicycle level of service at
intersections, roundabouts and other crossings}. Technical Report. Trafitec.
\url{https://trafitec.dk/wp-content/uploads/2023/07/158-Pedestrian-and-Bicycle-Level-of-Service-at-Intersections-Roundabouts-and-other-Crossings.pdf}.

\bibitem[Kazemzadeh et al.(2024)]{kazemzadeh2024forwhom}
Kazemzadeh, K., Afghari, A. P., \& Cherry, C. R. (2024). For whom is sharing
really scaring? Capturing unobserved heterogeneity in perceived comfort when
cycling in shared spaces. \emph{Transportation Research. Part F, Traffic
Psychology and Behaviour, 103}, 306--318.
\url{https://doi.org/10.1016/j.trf.2024.04.017}.

\bibitem[Keler et al.(2018)]{keler2018bicycle}
Keler, A., Kaths, J., Chucholowski, F., Chucholowski, M., Grigoropoulos, G.,
Spangler, M., Kaths, H., \& Busch, F. (2018). A bicycle simulator for
experiencing microscopic traffic flow simulation in urban environments. In
\emph{2018 21st international conference on intelligent transportation systems
(ITSC)} (pp. 3020--3023). Maui, HI: IEEE.
\url{https://ieeexplore.ieee.org/document/8569576/}.

\bibitem[Kwigizile et al.(2017)]{kwigizile2017real}
Kwigizile, V., Oh, J. S., Ikonomov, P., Hasan, R., Villalobos, C. G., Kurdi,
A. H., \& Shaw, A. (2017). \emph{Real time bicycle simulation study of
bicyclists' behaviors and their implication on safety}. Technical Report TRCLC
2015-03. Western Michigan University.
\url{https://rosap.ntl.bts.gov/view/dot/34885}.

\bibitem[Liu et al.(2012)]{liu2012response}
Liu, W. C., Jeng, M. C., Hwang, J. R., Doong, J. L., Lin, C. Y., \& Lai, C. H.
(2012). The response patterns of young bicyclists to a right-turning motorcycle:
A simulator study. \emph{Perceptual and Motor Skills, 115}, 385--402.
\url{https://doi.org/10.2466/22.25.27.PMS.115.5.385-402}. Publisher: SAGE
Publications Inc.

\bibitem[Mesimäki \& Luoma(2021)Mesimäki and Luoma]{mesimaki2021near}
Mesimäki, J., \& Luoma, J. (2021). Near accidents and collisions between
pedestrians and cyclists. \emph{European Transport Research Review, 13}, 38.
\url{https://doi.org/10.1186/s12544-021-00497-z}.
\url{https://etrr.springeropen.com/articles/10.1186/s12544-021-00497-z}.

\bibitem[Meuleners \& Fraser(2015)Meuleners and Fraser]{meuleners2015validation}
Meuleners, L., \& Fraser, M. (2015). A validation study of driving errors using
a driving simulator. \emph{Transportation Research. Part F, Traffic Psychology
and Behaviour, 29}, 14--21. \url{https://doi.org/10.1016/j.trf.2014.11.009}.
\url{https://www.sciencedirect.com/science/article/pii/S1369847814001764}.

\bibitem[Meuleners et al.(2023)]{meuleners2023improving}
Meuleners, L., Fraser, M., \& Roberts, P. (2023). Improving cycling safety
through infrastructure design: A bicycle simulator study. \emph{Transportation
Research Interdisciplinary Perspectives, 18}, Article 100768.
\url{https://doi.org/10.1016/j.trip.2023.100768}.
\url{https://www.sciencedirect.com/science/article/pii/S2590198223000155}.

\bibitem[Mohammadi et al.(2023)]{mohammadi2023how}
Mohammadi, A., Bianchi Piccinini, G., \& Dozza, M. (2023). How do cyclists
interact with motorized vehicles at unsignalized intersections? Modeling
cyclists' yielding behavior using naturalistic data. \emph{Accident Analysis and
Prevention, 190}, Article 107156.
\url{https://doi.org/10.1016/j.aap.2023.107156}.
\url{https://www.sciencedirect.com/science/article/pii/S0001457523002038}.

\bibitem[Mohammadi et al.(2024)]{mohammadi2024understanding}
Mohammadi, A., Bianchi Piccinini, G., \& Dozza, M. (2024). Understanding the
interaction between cyclists and motorized vehicles at unsignalized
intersections: Results from a cycling simulator study. \emph{Journal of Safety
Research, 90}, 306--318. \url{https://doi.org/10.1016/j.jsr.2024.05.007}.
\url{https://www.sciencedirect.com/science/article/pii/S0022437524000604}.

\bibitem[Nazemi et al.(2021)]{nazemi2021studying}
Nazemi, M., van Eggermond, M. A. B., Erath, A., Schaffner, D., Joos, M., \&
Axhausen, K. W. (2021). Studying bicyclists' perceived level of safety using a
bicycle simulator combined with immersive virtual reality. \emph{Accident
Analysis and Prevention, 151}, Article 105943.
\url{https://doi.org/10.1016/j.aap.2020.105943}.
\url{https://www.sciencedirect.com/science/article/pii/S0001457520317632}.

\bibitem[Negi \& Mitra(2020)Negi and Mitra]{negi2020fixation}
Negi, S., \& Mitra, R. (2020). Fixation duration and the learning process: An
eye tracking study with subtitled videos. \emph{Journal of Eye Movement
Research, 13}. \url{https://doi.org/10.16910/jemr.13.6.1}.
\url{https://www.ncbi.nlm.nih.gov/pmc/articles/PMC8012014/}.

\bibitem[Ng et al.(2017)]{ng2017cyclist}
Ng, A., Debnath, A. K., \& Heesch, K. C. (2017). Cyclist' safety perceptions of
cycling infrastructure at un-signalised intersections: Cross-sectional survey of
Queensland cyclists. \emph{Journal of Transport \& Health, 6}, 13--22.
\url{https://doi.org/10.1016/j.jth.2017.03.001}.
\url{https://www.sciencedirect.com/science/article/pii/S2214140516303188}.

\bibitem[Nikiforiadis \& Basbas(2019)Nikiforiadis and Basbas]{nikiforiadis2019can}
Nikiforiadis, A., \& Basbas, S. (2019). Can pedestrians and cyclists share the
same space? The case of a city with low cycling levels and experience.
\emph{Sustainable Cities and Society, 46}, Article 101453.
\url{https://doi.org/10.1016/j.scs.2019.101453}.

\bibitem[O'Hern et al.(2017)]{ohern2017validation}
O'Hern, S., Oxley, J., \& Stevenson, M. (2017). Validation of a bicycle
simulator for road safety research. \emph{Accident Analysis and Prevention,
100}, 53--58. \url{https://doi.org/10.1016/j.aap.2017.01.002}.
\url{https://www.sciencedirect.com/science/article/pii/S0001457517300027}.

\bibitem[Pashkevich et al.(2022)]{pashkevich2022visual}
Pashkevich, A., Burghardt, T. E., Puławska-Obiedowska, S., \& Šucha, M. (2022).
Visual attention and speeds of pedestrians, cyclists, and electric scooter
riders when using shared road -- a field eye tracker experiment. \emph{Case
Studies on Transport Policy, 10}, 549--558.
\url{https://doi.org/10.1016/j.cstp.2022.01.015}.

\bibitem[Poulos et al.(2017)]{poulos2017near}
Poulos, R. G., Hatfield, J., Rissel, C., Flack, L. K., Shaw, L., Grzebieta, R.,
\& McIntosh, A. S. (2017). Near miss experiences of transport and recreational
cyclists in New South Wales, Australia. Findings from a prospective cohort
study. \emph{Accident Analysis and Prevention, 101}, 143--153.
\url{https://doi.org/10.1016/j.aap.2017.01.020}.
\url{https://www.sciencedirect.com/science/article/pii/S0001457517300477}.

\bibitem[Pupil Core(2024)]{pupilcore2024}
Pupil Core (2024). \emph{Open source eye tracking platform}.
\url{https://pupil-labs.com/products/core}.

\bibitem[RoadRunner(2024)]{roadrunner2024}
RoadRunner (2024). \emph{RoadRunner}.
\url{https://de.mathworks.com/products/roadrunner.html}.

\bibitem[Scott-Deeter et al.(2023)]{scottdeeter2023assessing}
Scott-Deeter, L., Hurwitz, D., Russo, B., Smaglik, E., \& Kothuri, S. (2023).
Assessing the impact of three intersection treatments on bicyclist safety using
a bicycling simulator. \emph{Accident Analysis and Prevention, 179}, Article
106877. \url{https://doi.org/10.1016/j.aap.2022.106877}.
\url{https://www.sciencedirect.com/science/article/pii/S0001457522003128}.

\bibitem[Skoczyński(2021)]{skoczynski2021analysis}
Skoczyński, P. (2021). Analysis of solutions improving safety of cyclists in the
road traffic. \emph{Applied Sciences, 11}, 3771.
\url{https://doi.org/10.3390/app11093771}.
\url{https://www.mdpi.com/2076-3417/11/9/3771}.

\bibitem[Stülpnagel(2020)]{stulpnagel2020gaze}
Stülpnagel, R. (2020). Gaze behavior during urban cycling: Effects of subjective
risk perception and vista space properties. \emph{Transportation Research. Part
F, Traffic Psychology and Behaviour, 75}, 222--238.
\url{https://doi.org/10.1016/j.trf.2020.10.003}.

\bibitem[von Stülpnagel \& Rintelen(2024)von Stülpnagel and Rintelen]{vonstulpnagel2024matter}
von Stülpnagel, R., \& Rintelen, H. (2024). A matter of space and perspective --
cyclists', car drivers', and pedestrians' assumptions about subjective safety in
shared traffic situations. \emph{Transportation Research. Part A, Policy and
Practice, 179}, Article 103941.
\url{https://doi.org/10.1016/j.tra.2023.103941}.
\url{https://www.sciencedirect.com/science/article/pii/S0965856423003610}.

\bibitem[Suzuki et al.(2018)]{suzuki2018fundamental}
Suzuki, M., Sato, K., Hosoya, K., Miyanoue, K., \& Yai, T. (2018). \emph{A
fundamental study of validation of cycling simulator and designing safety
education programs}. \url{https://trid.trb.org/view/1495332}. Number: 18-02484.

\bibitem[Takanashi \& Murao(2018)Takanashi and Murao]{takanashi2018analysis}
Takanashi, H., \& Murao, I. (2018). Analysis of near-miss incidents between car
and bicycle at non-intersection. In \emph{Proceedings of international conference
on technology and social science 2018 (ICTSS 2018)}.
\url{http://conf.e-jikei.org/ICTSS2018/proceedings/materials/proc_files/IPS09\%5BHashimoto\%5D/ICTSS2018-IPS09-05/IPS09_05_takanashi.pdf}.

\bibitem[Velichkovsky et al.(2002)]{velichkovsky2002towards}
Velichkovsky, B. M., Rothert, A., Kopf, M., Dornhöfer, S. M., \& Joos, M.
(2002). Towards an express-diagnostics for level of processing and hazard
perception. \emph{Transportation Research. Part F, Traffic Psychology and
Behaviour, 5}, 145--156.
\url{https://doi.org/10.1016/S1369-8478(02)00013-X}.
\url{https://www.sciencedirect.com/science/article/pii/S136984780200013X}.

\bibitem[Wongwaiwit et al.(2011)]{wongwaiwit2011analysis}
Wongwaiwit, P., Michitsuji, Y., \& Raksincharoensak, Pongsathorn (2011).
Analysis on pedestrian and bicycle behavior in unsignalized intersection based
on near-miss incident database. In \emph{The proceedings of the transportation
and logistics conference 2011.20} (pp. 19--22).
\url{https://www.jstage.jst.go.jp/article/jsmetld/2011.20/0/2011.20_19/_article}.

\bibitem[Xu et al.(2023)]{xu2023assessing}
Xu, Z., Zheng, N., Logan, D. B., \& Vu, H. L. (2023). Assessing bicycle-vehicle
conflicts at urban intersections utilizing a VR integrated simulation approach.
\emph{Accident Analysis and Prevention, 191}, Article 107194.
\url{https://doi.org/10.1016/j.aap.2023.107194}.
\url{https://www.sciencedirect.com/science/article/pii/S0001457523002415}.

\bibitem[Zeuwts et al.(2023)]{zeuwts2023using}
Zeuwts, L. H. R. H., Vanhuele, R., Vansteenkiste, P., Deconinck, F. J. A., \&
Lenoir, M. (2023). Using an immersive virtual reality bicycle simulator to
evaluate hazard detection and anticipation of overt and covert traffic
situations in young bicyclists. \emph{Virtual Reality, 27}, 1507--1527.
\url{https://doi.org/10.1007/s10055-023-00746-7}.

\bibitem[Zhang et al.(2023)]{zhang2023space}
Zhang, C., Du, B., Zheng, Z., \& Shen, J. (2023). Space sharing between
pedestrians and micro-mobility vehicles: A systematic review.
\emph{Transportation Research. Part D, Transport and Environment, 116}, Article
103629. \url{https://doi.org/10.1016/j.trd.2023.103629}.
\url{https://www.sciencedirect.com/science/article/pii/S1361920923000263}.

\bibitem[Zhang et al.(2017)]{zhang2017study}
Zhang, R., Wu, J., Huang, L., \& You, F. (2017). Study of bicycle movements in
conflicts at mixed traffic unsignalized intersections. \emph{IEEE Access, 5},
10108--10117. \url{https://doi.org/10.1109/ACCESS.2017.2703816}.
\url{https://ieeexplore.ieee.org/document/7927409/}. Publisher: IEEE.

\end{thebibliography}
